\chapter{2017年山东卷（文）}

\section{单选题}

\begin{problem}
设集合 \(M=\{x\mid \lvert x-1\rvert<1\}\)，\(N=\{x\mid x<2\}\)，则 \(M\cap N=\)（\quad）.

\choices
    {\(\left(-1,1\right)\)}
    {\(\left(-1,2\right)\)}
    {\(\left(0,2\right)\)}
    {\(\left(1,2\right)\)}
\end{problem}

\begin{answer}
C
\end{answer}

\begin{solution}
由 \(\lvert x-1\rvert<1\) 得 \(-1<x-1<1\)，即 \(0<x<2\)，所以 \(M=(0,2)\).又 \(N=(-\infty,2)\)，故 \(M\cap N=(0,2)\)，选 C.
\end{solution}

\begin{problem}
已知 \(\i\) 是虚数单位，若复数 \(z\) 满足 \(z\i=1+\i\)，则 \(z^2=\)（\quad）.

\choices
    {\(-2\i\)}
    {\(2\i\)}
    {\(-2\)}
    {\(2\)}
\end{problem}

\begin{answer}
A
\end{answer}

\begin{solution}
由 \(\i^2=-1\)，将 \(z\i=1+\i\) 两边除以 \(\i\)，得 \(z=\dfrac{1+\i}{\i}=1-\i\).因此
\(z^2=(1-\i)^2=1-2\i+\i^2=-2\i\)，故选 A.
\end{solution}

\begin{problem}
已知 \(x\)，\(y\) 满足约束条件
\(\begin{cases}
x-2y+5\le0,\\
x+3\ge0,\\
y\le2,
\end{cases}\)
则 \(z=x+2y\) 的最大值是（\quad）.

\choices
    {\(-3\)}
    {\(-1\)}
    {\(1\)}
    {\(3\)}
\end{problem}

\begin{answer}
D
\end{answer}

\begin{solution}
约束条件分别等价于 \(y\geq\dfrac{x+5}{2}\)、\(x\geq-3\)、\(y\leq2\).可行域的顶点为 \((-3,1)\)、\((-3,2)\)、\((-1,2)\).在三个顶点处计算 \(z=x+2y\)，分别得到 \(-1\)、\(1\)、\(3\)，所以最大值为 \(3\)，选 D.
\end{solution}

\begin{problem}
已知 \(\cos x=\dfrac34\)，则 \(\cos2x=\)（\quad）.

\choices
    {\(-\dfrac14\)}
    {\(\dfrac14\)}
    {\(-\dfrac18\)}
    {\(\dfrac18\)}
\end{problem}

\begin{answer}
D
\end{answer}

\begin{solution}
由二倍角公式，
\(\cos2x=2\cos^2x-1=2\left(\dfrac34\right)^2-1=\dfrac{18}{16}-\dfrac{16}{16}=\dfrac18\).故选 D.
\end{solution}

\begin{problem}
已知命题 \(p\)：\(\exists x\in\R\)，\(x^2-x+1\ge0\)；命题 \(q\)：若 \(a^2<b^2\)，则 \(a<b\). 下列命题为真命题的是（\quad）.

\choices
    {\(p\land q\)}
    {\(p\land \lnot q\)}
    {\(\lnot p\land q\)}
    {\(\lnot p\land \lnot q\)}
\end{problem}

\begin{answer}
B
\end{answer}

\begin{solution}
因为 \(x^2-x+1=\left(x-\dfrac12\right)^2+\dfrac34>0\)，所以存在实数 \(x\) 使不等式成立，命题 \(p\) 为真.

命题 \(q\) 为假，例如取 \(a=1\)、\(b=-2\)，则 \(a^2<b^2\)，但 \(a<b\) 不成立.因此 \(p\land\lnot q\) 为真，选 B.
\end{solution}

\begin{problem}
执行右侧的程序框图，当输入的 \(x\) 的值为 \(4\) 时，输出的 \(y\) 的值为 \(2\)，则空白判断框中的条件可能为（\quad）.

\texfigure[max width=0.46\linewidth]{img_repaint/2017/shandong_liberal/q6_fig1.tex}

\choices
    {\(x>3\)}
    {\(x>4\)}
    {\(x\le4\)}
    {\(x\le5\)}
\end{problem}

\begin{answer}
B
\end{answer}

\begin{solution}
当输入 \(x=4\) 时，若判断框条件为真，程序沿“是”分支输出 \(y=x+2=6\)，与题设输出 \(2\) 不符.因此判断框条件在 \(x=4\) 时必须为假，且“否”分支给出 \(y=\log_2x=\log_2 4=2\).四个选项中只有 \(x>4\) 在 \(x=4\) 时为假，故选 B.
\end{solution}

\begin{problem}
函数 \(y=\sqrt3\sin2x+\cos2x\) 的最小正周期为（\quad）.

\choices
    {\(\dfrac{\pi}{2}\)}
    {\(\dfrac{2\pi}{3}\)}
    {\(\pi\)}
    {\(2\pi\)}
\end{problem}

\begin{answer}
C
\end{answer}

\begin{solution}
将函数合成为正弦型函数：\(y=\sqrt3\sin2x+\cos2x=2\sin\left(2x+\varphi\right)\)，其中 \(\varphi\) 为常数.正弦函数的最小正周期为 \(2\pi\)，自变量为 \(2x+\varphi\)，故原函数的最小正周期为 \(\dfrac{2\pi}{2}=\pi\)，选 C.
\end{solution}

\begin{problem}
如图所示的茎叶图记录了甲，乙两组各 5 名工人某日的产量数据（单位：件）. 若这两组数据的中位数相等，且平均值也相等，则 \(x\) 和 \(y\) 的值分别为（\quad）.

\begin{center}
\begin{tabular}{cc|c|ccc}
\multicolumn{2}{c|}{甲组} & \multicolumn{3}{c}{乙组} \\
\hline
 & \(6\) & \(5\) & \(9\) & & \\
\(2\) & \(5\) & \(6\) & \(1\) & \(7\) & \(y\) \\
\(x\) & \(4\) & \(7\) & \(8\) & &
\end{tabular}
\end{center}

\choices
    {\(3,5\)}
    {\(5,5\)}
    {\(3,7\)}
    {\(5,7\)}
\end{problem}

\begin{answer}
A
\end{answer}

\begin{solution}
茎叶图中茎为十位数字、叶为个位数字.甲组数据为 \(56\)，\(62\)，\(65\)，\(7x\)，\(74\)，其中位数为 \(65\).乙组数据为 \(59\)，\(61\)，\(67\)，\(6y\)，\(78\).要使乙组中位数也为 \(65\)，必须有 \(6y=65\)，所以 \(y=5\).

此时甲组数据的和为 \(56+62+65+(70+x)+74=327+x\)，乙组数据的和为 \(59+61+67+65+78=330\).平均值相等意味着两组数据的和相等，故 \(327+x=330\)，得 \(x=3\).因此 \((x,y)=(3,5)\)，选 A.
\end{solution}

\begin{problem}
设
\(f(x)=\begin{cases}
\sqrt{x},&0<x<1,\\
2(x-1),&x\ge1,
\end{cases}\)
若 \(f(a)=f(a+1)\)，则 \(f\!\left(\dfrac1a\right)=\)（\quad）.

\choices
    {\(2\)}
    {\(4\)}
    {\(6\)}
    {\(8\)}
\end{problem}

\begin{answer}
C
\end{answer}

\begin{solution}
由 \(f(a)=f(a+1)\) 可知 \(a>0\).又因 \(a+1\geq1\)，所以 \(f(a+1)=2a\).

当 \(0<a<1\) 时，\(f(a)=\sqrt a\)，于是 \(\sqrt a=2a\).由于 \(a>0\)，除以 \(\sqrt a\) 得 \(1=2\sqrt a\)，所以 \(a=\dfrac14\)，从而
\(f\left(\dfrac1a\right)=f(4)=2(4-1)=6\).

当 \(a\geq1\) 时，\(f(a)=2(a-1)\)，等式变为 \(2(a-1)=2a\)，矛盾.因此唯一可能的值为 \(6\)，选 C.
\end{solution}

\begin{problem}
若函数 \(\e^x f(x)\)（\(\e=2.71828\ldots\) 是自然对数的底数）在 \(f(x)\) 的定义域上单调递增，则称函数 \(f(x)\) 具有 \(M\) 性质. 下列函数中具有 \(M\) 性质的是（\quad）.

\choices
    {\(f(x)=2^{-x}\)}
    {\(f(x)=x^2\)}
    {\(f(x)=3^{-x}\)}
    {\(f(x)=\cos x\)}
\end{problem}

\begin{answer}
A
\end{answer}

\begin{solution}
对选项 A，\(\e^xf(x)=\e^x2^{-x}=\left(\dfrac{\e}{2}\right)^x\).因为 \(\e>2\)，所以 \(\dfrac{\e}{2}>1\)，该函数在 \(\R\) 上单调递增，故 \(f(x)=2^{-x}\) 具有 \(M\) 性质.

选项 B 中 \((\e^xx^2)'=\e^x(x^2+2x)\)，导数在 \((-2,0)\) 与其余区间符号不同；选项 C 中 \(\left(\dfrac{\e}{3}\right)^x\) 单调递减；选项 D 中 \((\e^x\cos x)'=\e^x(\cos x-\sin x)\) 的符号发生变化.因此选 A.
\end{solution}

\section{填空题}

\begin{problem}
已知向量 \(\bs{a}=(2,6)\)，\(\bs{b}=(-1,\lambda)\)，若 \(\bs{a}\parallel\bs{b}\)，则 \(\lambda=\fillinblank{}\).
\end{problem}

\begin{answer}
\(-3\)
\end{answer}

\begin{solution}
两个非零向量平行，当且仅当其坐标行列式为零.因此
\(\begin{vmatrix}2&6\\-1&\lambda\end{vmatrix}=2\lambda+6=0\)，解得 \(\lambda=-3\).
\end{solution}

\begin{problem}
若直线 \(\dfrac{x}{a}+\dfrac{y}{b}=1\)（\(a>0\)，\(b>0\)）过点 \((1,2)\)，则 \(2a+b\) 的最小值为\fillinblank{}.
\end{problem}

\begin{answer}
8
\end{answer}

\begin{solution}
直线过点 \((1,2)\)，所以
\(\dfrac1a+\dfrac2b=1\)，其中 \(a>0\)，\(b>0\).由柯西不等式，
\[
(2a+b)\left(\dfrac1a+\dfrac2b\right)
\geq\left(\sqrt{2a\cdot\dfrac1a}+\sqrt{b\cdot\dfrac2b}\right)^2=8
\]
由于 \(\dfrac1a+\dfrac2b=1\)，故 \(2a+b\geq8\).当且仅当 \(\dfrac{1/a}{2a}=\dfrac{2/b}{b}\)，即 \(b=2a\) 时取等号；代入约束得 \(a=2\)，\(b=4\).因此最小值为 \(8\).
\end{solution}

\begin{problem}
由一个长方体和两个 \(\dfrac14\) 圆柱构成的几何体的三视图如右图，则该几何体的体积为\fillinblank{}.

\bitmapfigure[max width=0.45\linewidth]{img/2017/shandong_liberal/q12_fig1.png}
\end{problem}

\begin{answer}
\(2+\dfrac{\pi}{2}\)
\end{answer}

\begin{solution}
由三视图可知，长方体的长、宽、高分别为 \(2\)，\(1\)，\(1\)，故其体积为 \(2\).两端各为一个半径为 \(1\)、高为 \(1\) 的四分之一圆柱，每个体积为 \(\dfrac14\pi\cdot1^2\cdot1=\dfrac\pi4\).因此该几何体的体积为
\(2+2\cdot\dfrac\pi4=2+\dfrac\pi2\).
\end{solution}

\begin{problem}
已知 \(f(x)\) 是定义在 \(\R\) 上的偶函数，且 \(f(x+4)=f(x-2)\). 若当 \(x\in\left[-3,0\right]\) 时，\(f(x)=6^{-x}\)，则 \(f(919)=\fillinblank{}\).
\end{problem}

\begin{answer}
6
\end{answer}

\begin{solution}
由 \(f(x+4)=f(x-2)\)，令 \(t=x-2\)，得 \(f(t+6)=f(t)\)，所以 \(6\) 是函数的一个周期.因 \(919=153\times6+1\)，有 \(f(919)=f(1)\).

函数为偶函数，所以 \(f(1)=f(-1)\).又 \(-1\in[-3,0]\)，由已知条件得 \(f(-1)=6^{-(-1)}=6\).故 \(f(919)=6\).
\end{solution}

\begin{problem}
在平面直角坐标系 \(xOy\) 中，双曲线 \(\dfrac{x^2}{a^2}-\dfrac{y^2}{b^2}=1\)（\(a>0\)，\(b>0\)）的右支与焦点为 \(F\) 的抛物线 \(x^2=2py\)（\(p>0\)）交于 \(A\)，\(B\) 两点，若 \(\lvert AF\rvert+\lvert BF\rvert=4\lvert OF\rvert\)，则该双曲线的渐近线方程为\fillinblank{}.
\end{problem}

\begin{answer}
\(y=\pm\dfrac{\sqrt{2}}{2}x\)
\end{answer}

\begin{solution}
抛物线 \(x^2=2py\) 的焦点为 \(F\left(0,\dfrac p2\right)\)，所以 \(\lvert OF\rvert=\dfrac p2\).对抛物线上任意点 \((x,y)\)，由抛物线的定义，点到焦点的距离等于点到准线 \(y=-\dfrac p2\) 的距离，即 \(\lvert (x,y)F\rvert=y+\dfrac p2\).

设交点 \(A\)，\(B\) 的纵坐标分别为 \(y_A\)，\(y_B\)，则题设给出
\[
\lvert AF\rvert+\lvert BF\rvert=y_A+y_B+p=4\lvert OF\rvert=2p
\]
所以 \(y_A+y_B=p\).将抛物线方程代入双曲线方程，得
\(\dfrac{2py}{a^2}-\dfrac{y^2}{b^2}=1\)，即
\(\dfrac{y^2}{b^2}-\dfrac{2p}{a^2}y+1=0\).由韦达定理，两个交点纵坐标之和为 \(\dfrac{2pb^2}{a^2}\)，故
\(\dfrac{2pb^2}{a^2}=p\)，从而 \(\dfrac{b^2}{a^2}=\dfrac12\).

双曲线的渐近线为 \(y=\pm\dfrac ba x\)，因此方程为 \(y=\pm\dfrac{\sqrt2}{2}x\).
\end{solution}

\section{解答题}

\begin{problem}
某旅游爱好者计划从 3 个亚洲国家 \(A_1\)，\(A_2\)，\(A_3\) 和 3 个欧洲国家 \(B_1\)，\(B_2\)，\(B_3\) 中选择 2 个国家去旅游.

\begin{enumerate}
\item 若从这 6 个国家中任选 2 个，求这 2 个国家都是亚洲国家的概率；
\item 若从亚洲国家和欧洲国家中各任选 1 个，求这 2 个国家包括 \(A_1\) 但不包括 \(B_1\) 的概率.
\end{enumerate}
\end{problem}

\begin{solution}
\begin{enumerate}
\item 从 6 个国家中任选 2 个国家，共有 \(\mathrm C_6^2=15\) 个等可能结果.两国都是亚洲国家的选法有 \(\mathrm C_3^2=3\) 个，所以所求概率为
\[
\dfrac{\mathrm C_3^2}{\mathrm C_6^2}=\dfrac{3}{15}=\dfrac15
\]

\item 从亚洲国家和欧洲国家中各选 1 个，共有 \(3\times3=9\) 个等可能结果.若包括 \(A_1\) 且不包括 \(B_1\)，亚洲国家必须选 \(A_1\)，欧洲国家可选 \(B_2\) 或 \(B_3\)，共有 2 个结果.因此所求概率为
\[
\dfrac{2}{9}
\]
\end{enumerate}
\end{solution}

\begin{problem}
在 \(\triangle ABC\) 中，角 \(A\)，\(B\)，\(C\) 的对边分别为 \(a\)，\(b\)，\(c\)，已知 \(b=3\)，\(\overrightarrow{AB}\cdot\overrightarrow{AC}=-6\)，\(S_{\triangle ABC}=3\)，求 \(A\) 和 \(a\).
\end{problem}

\begin{solution}
设 \(c=\lvert AB\rvert\).由向量点积公式，
\[
\overrightarrow{AB}\cdot\overrightarrow{AC}=\lvert AB\rvert\lvert AC\rvert\cos A=3c\cos A=-6
\]
所以 \(c\cos A=-2\).又由三角形面积公式，
\[
S_{\triangle ABC}=\dfrac12\lvert AB\rvert\lvert AC\rvert\sin A=\dfrac32c\sin A=3
\]
所以 \(c\sin A=2\).

因为 \(0<A<\pi\)，有 \(\sin A>0\)，且由上式 \(\cos A<0\)，所以 \(A\) 为钝角.将 \(c\cos A=-2\)、\(c\sin A=2\) 平方相加，得
\(c^2=8\)，故 \(c=2\sqrt2\).于是 \(\cos A=-\dfrac{1}{\sqrt2}\)、\(\sin A=\dfrac{1}{\sqrt2}\)，从而 \(A=\dfrac{3\pi}{4}\).

由余弦定理，
\[
a^2=b^2+c^2-2bc\cos A=3^2+(2\sqrt2)^2-2\cdot3\cdot2\sqrt2\cdot\left(-\dfrac{\sqrt2}{2}\right)=29
\]
因 \(a>0\)，得 \(a=\sqrt{29}\).
\end{solution}

\begin{problem}
由四棱柱 \(ABCD-A_1B_1C_1D_1\) 截去三棱锥 \(C_1-B_1CD_1\) 后得到的几何体如图所示，四边形 \(ABCD\) 为正方形，\(O\) 为 \(AC\) 与 \(BD\) 的交点，\(E\) 为 \(AD\) 的中点，\(A_1E\perp\) 平面 \(ABCD\).

\begin{enumerate}
\item 证明：\(A_1O\parallel\) 平面 \(B_1CD_1\)；
\item 设 \(M\) 是 \(OD\) 的中点，证明：平面 \(A_1EM\perp\) 平面 \(B_1CD_1\).
\end{enumerate}

\bitmapfigure[max width=0.58\linewidth]{img/2017/shandong_liberal/q18_fig1.png}
\end{problem}

\begin{solution}
设正方形边长为 \(1\)，以平面 \(ABCD\) 为 \(xOy\) 平面，并取
\(A=(0,0,0)\)，\(B=(1,0,0)\)，\(C=(1,1,0)\)，\(D=(0,1,0)\).
因为 \(E\) 为 \(AD\) 的中点且 \(A_1E\perp\) 平面 \(ABCD\)，可设 \(A_1=(0,\frac12,h)\)，其中 \(h>0\).四棱柱的上底面由同一个平移向量得到，故
\(B_1=(1,\tfrac12,h)\)，\(D_1=(0,\tfrac32,h)\).
又 \(O=(\frac12,\frac12,0)\)，且 \(M\) 为 \(OD\) 的中点，所以 \(M=(\frac14,\frac34,0)\).

（1）有
\(\overrightarrow{A_1O}=(\tfrac12,0,-h)\)，\(\overrightarrow{B_1C}=(0,\tfrac12,-h)\)，\(\overrightarrow{B_1D_1}=(-1,1,0)\).
并且
\[
\overrightarrow{A_1O}=\overrightarrow{B_1C}-\tfrac12\overrightarrow{B_1D_1}
\]
所以直线 \(A_1O\) 的方向向量是平面 \(B_1CD_1\) 内两个相交方向向量的线性组合，即直线 \(A_1O\) 的方向平行于该平面.平面 \(B_1CD_1\) 的法向量可取 \((h,h,\tfrac12)\)，而 \(\overrightarrow{B_1A_1}=(-1,0,0)\) 与该法向量的内积为 \(-h\ne0\)，故 \(A_1\) 不在该平面内，从而直线 \(A_1O\parallel\) 平面 \(B_1CD_1\).

（2）平面 \(A_1EM\) 内两条相交向量为
\(\overrightarrow{EA_1}=(0,0,h)\)，\(\overrightarrow{EM}=(\tfrac14,\tfrac14,0)\)
其法向量可取
\[
\bs{n}_1=\overrightarrow{EA_1}\times\overrightarrow{EM}=(-\tfrac h4,\tfrac h4,0)
\]
平面 \(B_1CD_1\) 的一个法向量为
\[
\bs{n}_2=\overrightarrow{B_1C}\times\overrightarrow{B_1D_1}=(h,h,\tfrac12)
\]
于是 \(\bs{n}_1\cdot\bs{n}_2=-\dfrac{h^2}{4}+\dfrac{h^2}{4}=0\).两平面的法向量垂直，所以平面 \(A_1EM\perp\) 平面 \(B_1CD_1\).
\end{solution}

\begin{problem}
已知 \(\{a_n\}\) 是各项均为正数的等比数列，且 \(a_1+a_2=6\)，\(a_1a_2=a_3\).

\begin{enumerate}
\item 求数列 \(\{a_n\}\) 的通项公式；
\item \(\{b_n\}\) 为各项非零的等差数列，其前 \(n\) 项和为 \(S_n\)，已知 \(S_{2n+1}=b_nb_{n+1}\)，求数列 \(\left\{\dfrac{b_n}{a_n}\right\}\) 的前 \(n\) 项和 \(T_n\).
\end{enumerate}
\end{problem}

\begin{solution}
（1）设等比数列 \(\{a_n\}\) 的公比为 \(q>0\).由 \(a_1a_2=a_3\)，得
\[
a_1\cdot a_1q=a_1q^2
\]
因各项为正数，可约去 \(a_1q\)，得到 \(a_1=q\).再由 \(a_1+a_2=6\)，得 \(q+q^2=6\)，即 \((q-2)(q+3)=0\).由于 \(q>0\)，所以 \(q=2\)，且 \(a_1=2\).因此
\[
a_n=a_1q^{n-1}=2\cdot2^{n-1}=2^n
\]

（2）等差数列的前 \(2n+1\) 项和等于中间项 \(b_{n+1}\) 乘以项数，因此 \(S_{2n+1}=(2n+1)b_{n+1}\).由已知等式，
\[
(2n+1)b_{n+1}=b_nb_{n+1}
\]
由于各项非零，特别是 \(b_{n+1}\ne0\)，故 \(b_n=2n+1\).于是
\[
T_n=\sum_{k=1}^n\dfrac{b_k}{a_k}=\sum_{k=1}^n\dfrac{2k+1}{2^k}
\]
注意到
\[
\dfrac{2k+1}{2^k}=\dfrac{4k+6}{2^k}-\dfrac{4(k+1)+6}{2^{k+1}}
\]
所以裂项相消得
\[
T_n=\dfrac{10}{2}-\dfrac{4n+10}{2^{n+1}}=5-\dfrac{2n+5}{2^n}
\]
\end{solution}

\begin{problem}
已知函数 \(f(x)=\dfrac13x^3-\dfrac a2x^2\)，\(a\in\R\).

\begin{enumerate}
\item 当 \(a=2\) 时，求曲线 \(y=f(x)\) 在点 \((3,f(3))\) 处的切线方程；
\item 设函数 \(g(x)=f(x)+(x-a)\cos x-\sin x\)，讨论 \(g(x)\) 的单调性并判断有无极值，有极值时求出极值.
\end{enumerate}
\end{problem}

\begin{solution}
（1）当 \(a=2\) 时，\(f(3)=\dfrac13\cdot3^3-\dfrac22\cdot3^2=0\)，且
\(f'(x)=x^2-ax\)，所以 \(f'(3)=9-6=3\).曲线在 \((3,0)\) 处的切线为
\[
y-0=3(x-3)
\]
即 \(3x-y-9=0\).

（2）由
\[
g'(x)=f'(x)+\cos x-(x-a)\sin x-\cos x=(x-a)(x-\sin x)
\]
令 \(h(x)=x-\sin x\)，则 \(h'(x)=1-\cos x\geq0\)，且 \(h(0)=0\)，所以 \(h(x)<0\)（\(x<0\)），\(h(x)>0\)（\(x>0\)）.因此 \(g'(x)\) 的符号由 \((x-a)x\) 决定，下面分情况讨论.

当 \(a<0\) 时，\(g(x)\) 在 \(\left(-\infty,a\right)\)、\(\left(a,0\right)\)、\(\left(0,+\infty\right)\) 上分别单调递增、递减、递增.在 \(x=a\) 处由增变减，故有极大值
\[
g(a)=f(a)-\sin a=-\dfrac16a^3-\sin a;
\]
在 \(x=0\) 处由减变增，故有极小值 \(g(0)=-a\).

当 \(a=0\) 时，\(g'(x)=x(x-\sin x)\geq0\)，且除 \(x=0\) 外为正，所以 \(g(x)\) 在 \(\R\) 上单调递增，无极值.

当 \(a>0\) 时，\(g(x)\) 在 \(\left(-\infty,0\right)\)、\(\left(0,a\right)\)、\(\left(a,+\infty\right)\) 上分别单调递增、递减、递增.在 \(x=0\) 处有极大值 \(g(0)=-a\)，在 \(x=a\) 处有极小值
\[
g(a)=-\dfrac16a^3-\sin a
\]
\end{solution}

\begin{problem}
在平面直角坐标系 \(xOy\) 中，已知椭圆 \(C:\dfrac{x^2}{a^2}+\dfrac{y^2}{b^2}=1\)（\(a>b>0\)）的离心率为 \(\dfrac{\sqrt2}{2}\)，椭圆 \(C\) 截直线 \(y=1\) 所得线段的长度为 \(2\sqrt2\).

\begin{enumerate}
\item 求椭圆 \(C\) 的方程；
\item 动直线 \(l:y=kx+m\)（\(m\ne0\)）交椭圆 \(C\) 于 \(A\)，\(B\) 两点，交 \(y\) 轴于点 \(M\). 点 \(N\) 是 \(M\) 关于 \(O\) 的对称点，\(\odot N\) 的半径为 \(\lvert NO\rvert\). 设 \(D\) 为 \(AB\) 的中点，\(DE\)，\(DF\) 与 \(\odot N\) 分别相切于点 \(E\)，\(F\)，求 \(\angle EDF\) 的最小值.
\end{enumerate}

\bitmapfigure[max width=0.55\linewidth]{img/2017/shandong_liberal/q21_fig1.png}
\end{problem}

\begin{solution}
（1）设椭圆的半焦距为 \(c\)，则
\(\dfrac ca=\dfrac{\sqrt2}{2}\)，\(c^2=a^2-b^2\).
所以 \(c^2=\dfrac{a^2}{2}\)，即 \(b^2=\dfrac{a^2}{2}\).椭圆与直线 \(y=1\) 的交点横坐标满足
\(\dfrac{x^2}{a^2}+\dfrac1{b^2}=1\).题设截得线段长度为 \(2\sqrt2\)，故两个交点横坐标为 \(\pm\sqrt2\)，从而
\[
\dfrac2{a^2}+\dfrac1{b^2}=1
\]
代入 \(b^2=\dfrac{a^2}{2}\)，得 \(\dfrac4{a^2}=1\)，所以 \(a^2=4\)，\(b^2=2\).椭圆方程为
\[
\dfrac{x^2}{4}+\dfrac{y^2}{2}=1
\]

（2）由（1）知椭圆为 \(\dfrac{x^2}{4}+\dfrac{y^2}{2}=1\).直线 \(l:y=kx+m\) 与椭圆联立，得到关于交点横坐标的方程
\[
\left(1+2k^2\right)x^2+4kmx+2m^2-4=0
\]
设 \(Q=1+2k^2\).由韦达定理，弦 \(AB\) 中点 \(D\) 的坐标为
\[
D=\left(-\dfrac{2km}{Q},\dfrac{m}{Q}\right)
\]
又 \(M=(0,m)\)，\(N=(0,-m)\)，故
\[
DN^2=\dfrac{4m^2\left(k^4+3k^2+1\right)}{\left(1+2k^2\right)^2}
\]
令 \(t=k^2\geq0\)，则
\[
\dfrac{DN^2}{m^2}=4\cdot\dfrac{t^2+3t+1}{(2t+1)^2}\leq4
\]
因为 \((2t+1)^2-(t^2+3t+1)=3t^2+t\geq0\).等号当且仅当 \(t=0\)，即 \(k=0\).

设 \(\theta=\angle EDF\)，圆的半径为 \(NO=|m|\).两条切线长相等，故 \(DN\) 平分 \(\angle EDF\).在直角三角形 \(DNE\) 中，
\[
\sin\dfrac\theta2=\dfrac{NE}{DN}=\dfrac{|m|}{DN}\geq\dfrac12
\]
由于 \(0<\theta<\pi\)，得 \(\theta\geq2\arcsin\dfrac12=\dfrac\pi3\).当 \(k=0\) 且取 \(0<|m|<\sqrt2\) 时，直线为水平割线并确实与椭圆交于两点，此时等号成立.因此 \(\angle EDF\) 的最小值为 \(\dfrac\pi3\).
\end{solution}
